\section*{Chapter \ref{ch:g8:the-oceans} --- The Oceans}

\begin{solution}{exo:g8:the-oceans:1}
\omterm{def:g8:the-oceans:seafloor}{Continental shelf}, \omterm{def:g8:the-oceans:seafloor}{continental slope}, \omterm{def:g8:the-oceans:seafloor}{abyssal plain}, \omterm{def:g8:the-oceans:seafloor}{ocean ridge}.
\end{solution}

\begin{solution}{exo:g8:the-oceans:2}
(a) False: about \qty{3700}{m}, that is \qty{3.7}{km}. (b) True. (c)
False: evaporation removes water and leaves the salts, so the water gets
saltier. (d) True.
\end{solution}

\begin{solution}{exo:g8:the-oceans:3}
(a) An \omterm{def:g8:the-oceans:seafloor}{ocean ridge}. (b) The \omterm{def:g8:the-oceans:salinity}{salinity}. (c) A \omterm{def:g8:the-oceans:current}{gyre}. (d) An \omterm{def:g8:the-oceans:seafloor}{ocean trench}.
\end{solution}

\begin{solution}{exo:g8:the-oceans:4}
(a) Anticlockwise. (b) Warm: the Gulf Stream; cold: the Canary Current.
(c) The \omterm{def:g8:winds-global-circulation:global-winds}{trade winds}.
\end{solution}

\begin{solution}{exo:g8:the-oceans:5}
(a) $35 \div 1000 \times 100 = 3.5~\%$. (b) $20 \times 35 = \qty{700}{g}$.
(c) $1000 \div 35 \approx 28.6$: about \qty{29}{kg}.
\end{solution}

\begin{solution}{exo:g8:the-oceans:6}
(a) $1500 \times 6 = \qty{9000}{m}$. (b) The sound goes down to the bottom
and comes back up: it covers the depth twice. (c) $9000 \div 2 =
\qty{4500}{m}$. (d) An \omterm{def:g8:the-oceans:seafloor}{abyssal plain} (a trench would be deeper, a ridge or
a shelf shallower).
\end{solution}

\begin{solution}{exo:g8:the-oceans:7}
(a) Between about \qty{3500}{m} and \qty{5000}{m}. (b) About \qty{4400}{km} along the profile, about \qty{2300}{m} deep. (c) About
$4500 - 2300 = \qty{2200}{m}$ (between 1200 and \qty{2700}{m}, depending
on the \omterm{def:g1:landscapes-around-us:landform}{plain} chosen). (d) The \omterm{def:g2:water-on-earth:ocean}{ocean} is a few kilometres deep but thousands
of kilometres wide: drawn at the same scale, the whole profile would be
thinner than a line.
\end{solution}

\begin{solution}{exo:g8:the-oceans:8}
\omterm{def:g2:water-on-earth:ocean}{Sea} A is saltier: the Sun and the dry air cause strong
\emph{evaporation}, the water leaves, the \emph{salts stay behind}, and
almost no \omterm{def:g4:clouds-rain-snow:precipitation}{rain} or \omterm{def:g1:landscapes-around-us:water}{river} brings \emph{\omterm{def:g2:water-on-earth:salt-fresh}{fresh water}} back. \omterm{def:g2:water-on-earth:ocean}{Sea} B is less
salty: its \omterm{def:g1:landscapes-around-us:water}{rivers} and its \omterm{def:g4:clouds-rain-snow:precipitation}{rain} bring a lot of \emph{\omterm{def:g2:water-on-earth:salt-fresh}{fresh water}}, which
dilutes the salts.
\end{solution}

\begin{solution}{exo:g8:the-oceans:9}
(a) The Gulf Stream, towards the north along the \omterm{def:g1:landscapes-around-us:water}{coast}, then towards the
north-east. (b) Towards the east (north-east), with the North Atlantic
Current. (c) Along the \omterm{def:g1:landscapes-around-us:water}{coast} of north-west Africa (and Spain and
Portugal), with the Canary Current. (d) Cold.
\end{solution}

\begin{solution}{exo:g8:the-oceans:10}
(a) $6.4 \times 24 = 153.6$: about \qty{154}{km}. (b) $153.6 \times 7
\approx \qty{1075}{km}$. (c) The Gulf Stream is a little faster than a
walker, but it never stops, day and night.
\end{solution}

\begin{solution}{exo:g8:the-oceans:11}
(a) \qty{3.7}{km} and \qty{11}{km}. (b) $11000 \div 3700 \approx 3.0$
times. (c) The sound travels $2 \times 11000 = \qty{22000}{m}$; $22000
\div 1500 \approx 14.7$: about \qty{15}{s}. (d) $1000 = 10^3$ years.
\end{solution}

\begin{solution}{exo:g8:the-oceans:12}
(a) Near the \omterm{def:g3:continents-and-oceans:equator}{equator}, the \omterm{def:g8:winds-global-circulation:global-winds}{trade winds} push warm tropical water towards the
west. At the western edge of the Atlantic it turns north with the Gulf
Stream, a warm current, along Florida and North America. The \omterm{def:g8:winds-global-circulation:global-winds}{westerlies}
then push it across the Atlantic towards the north-east, as far as the
\omterm{def:g2:water-on-earth:ocean}{seas} off Norway. (b) There it gives its heat to the cold air: it cools,
and it becomes saltier when \omterm{def:g2:water-on-earth:ocean}{sea} ice forms, because the salt is left in
the water. Cold, salty water is dense, so it sinks to the bottom. (c) It
creeps south along the \omterm{def:g2:water-on-earth:ocean}{ocean} floor, towards Antarctica, then into the
Indian and Pacific \omterm{def:g2:water-on-earth:ocean}{Oceans}, where it slowly rises again.
\end{solution}

\begin{solution}{pb:g8:the-oceans:1}
\textbf{1.} In the Atlantic \omterm{def:g2:water-on-earth:ocean}{Ocean}, in the \omterm{def:g3:continents-and-oceans:equator}{northern hemisphere}.

\textbf{2.} Clockwise; the buoy is carried north at first.

\textbf{3.} Leg A (Gulf Stream) is warm: it comes from the \omterm{def:g5:weather-and-climate:tropics}{tropics}. Leg C
(Canary Current) is cold: it comes from higher \omterm{def:g6:latitude-longitude:latitude}{latitudes}, towards the
\omterm{def:g3:continents-and-oceans:equator}{equator}.

\textbf{4.} Leg B: the \omterm{def:g8:winds-global-circulation:global-winds}{westerlies}. Leg D: the \omterm{def:g8:winds-global-circulation:global-winds}{trade winds}.

\textbf{5.} $3800 \div 100 = 38$ days.

\textbf{6.} B: $2100 \div 20 = 105$ days. C: $3100 \div 20 = 155$ days.
D: $6100 \div 20 = 305$ days.

\textbf{7.} $3800 + 2100 + 3100 + 6100 = \qty{15100}{km}$.

\textbf{8.} The Gulf Stream is the narrow, fast current on the western
side of the \omterm{def:g8:the-oceans:current}{gyre}; the currents on the other sides are broad and slow.

\textbf{9.} $38 + 105 + 155 + 305 = 603$ days.

\textbf{10.} $603 \div 30 = 20.1$: about \textbf{20 months}.

\textbf{11.} After legs A, B and C, $38 + 105 + 155 = 298$ days have
passed; 365 days come during leg D (the North Equatorial Current), $365 -
298 = 67$ days into it.

\textbf{12.} $603 - 365 = 238$ more days, that is $238 \div 30 \approx 8$
more months.
\end{solution}
